\chapter*{ملخص المشروع}
\addcontentsline{toc}{chapter}{ملخص المشروع باللغة العربية}

تتطلب النهضة العمرانية المتسارعة والمشاريع الكبرى والمبادرات البيئية والزراعية في المملكة العربية السعودية تحقيقاً لمستهدفات رؤية المملكة 2030 وجود منظومات تقنية مؤتمتة وقابلة للتوسع للرصد البيئي المستمر. غير أن الكشف الآلي عن التغيرات المكانية عبر الأقمار الاصطناعية في البيئات شديدة الجفاف يواجه تحديات علمية وبيئية بالغة التعقيد؛ حيث تؤدي التغيرات الفصلية في زوايا الإشعاع الشمسي، وتفاوت رطوبة التربة والسبخات الساحلية، والتشابه الطيفي الحاد بين الخرسانة الحضرية والصخور الجيرية الصحراوية إلى ارتفاع هائل في معدلات الإنذارات الكاذبة لدى خوارزميات الرؤية الحاسوبية التقليدية ونماذج الأساس العالمية المستقلة.

يقدم هذا المشروع التطبيقي منصة \textbf{TerraShift AI} (إعداد الطالب: \textbf{Omar Ayyash}، الرقم الجامعي: \textbf{S22107722})، وهي منصة جغرافية مكانية متكاملة ومصممة هندسياً للتعرف الآلي، وتحديد الموقع، وقياس، وتحليل التغيرات الواقعية على سطح الأرض بالاعتماد على صور الأقمار الاصطناعية المجانية من برنامج كوبرنيكوس الأوروبي (Sentinel-2 Level-2A) ومتجهات التضمين السنوية لنماذج الأساس العالمية (\texttt{GOOGLE/SATELLITE\_EMBEDDING/V1/ANNUAL} AlphaEarth 64-D) واحتمالات الغطاء الأرضي (\texttt{GOOGLE/DYNAMICWORLD/V1}) المستخرجة مباشرة عبر منصة Google Earth Engine الموثقة (\texttt{project='terrashift'}, \texttt{allow\_synthetic=False}). وللإجابة علمياً على التساؤل البحثي المحوري للمشروع عبر تقييم خالٍ تماماً من البيانات الاصطناعية، أجرينا دراسة مقارنة شاملة عبر أربعة ممرات مرجعية سعودية بمساحة إجمالية تبلغ 91.19 كم² وتضم 19 حدث تغير مرجعي و4 مناطق ضبط سلبية غير مضطربة: (1) ممر روشن سدرة بالرياض، (2) ممر شمال أبحر بجدة، (3) حوض وادي السرحان الزراعي بالجوف، و(4) ممر الخرج الصناعي واللوجستي (ممر اختبار مستقل ومحجوب لقياس التعميم الجغرافي).

منهجياً، قمنا بصياغة بروتوكول مطابقة أحادي صارم يراعي دقة التوطين المكاني (Localization-Aware 1-to-1 Matching: $\text{PoC} \ge 0.15, \text{IoU} \ge 0.05$) ويرفض المضلعات الضخمة غير الدقيقة، مع القص الهندسي الدقيق للمضلعات ضمن حدود منطقة الاهتمام (WGS84 Ellipsoidal Geodesic Clipping) وإعلان النتائج المجمعة (Pooled/Micro) والمتوسطة عبر الممرات (Corridor Macro-Averaged) جنباً إلى جنب عبر سبعة نماذج: (1) النموذج الهجين \texttt{alphaearth\_hybrid} (TerraShift v6.0)، (2) النموذج السيامي المضبوط المدعم بالحواجز الفيزيائية \texttt{siamese\_tuned\_filtered} (TerraShift v5.1)، (3) نموذج \texttt{alphaearth\_cosine} المستقل، (4) نموذج الفيزياء الطيفية البحتة \texttt{multispectral\_physical}، (5) الشبكة السيامية غير المضبوطة \texttt{siamese\_untuned}، (6) نموذج \texttt{dynamic\_world\_lulc} المستقل، و(7) نموذج \texttt{alphaearth\_dynamic\_world}.

أثبتت النتائج التجريبية الموثقة أن النموذج الهجين \textbf{TerraShift v6.0 (\texttt{alphaearth\_hybrid})} حقق المرتبة الأولى بمقياس F1 المجمع (Pooled F1) بلغ \textbf{79.2\%} ومقياس F1 المتوسط (Macro F1) بلغ \textbf{82.5\%}، مع دقة مجمعة \textbf{65.5\%} (ودقة متوسطة \textbf{70.1\%}) واستدعاء كامل \textbf{100.0\%} (19/19)، وتقليص الإنذارات الكاذبة بنسبة \textbf{91.2\%} (من 114 إلى 10 إنذارات فقط) مع صفر إنذارات في مناطق الضبط السلبية. كما حقق نموذج \textbf{TerraShift v5.1 (\texttt{siamese\_tuned\_filtered})} المرتبة الثانية (\textbf{53.5\%} Pooled F1 / \textbf{64.1\%} Macro F1، واستدعاء 100\%، وتقليص الإنذارات الكاذبة بنسبة 71.1\%) مع ميزة الرصد التكتيكي كل 5 أيام دون انتظار المتجهات السنوية. وفي المقابل، حقق نموذج AlphaEarth المستقل \textbf{27.0\%} Pooled F1 فقط بسبب إطلاقه 90 إنذاراً كاذباً (منها 23 إنذاراً في مناطق الضبط السلبية الصحراوية والسبخات)، بينما حقق نموذج Dynamic World المستقل \textbf{19.7\%} Pooled F1 فقط (استدعاء 31.6\%).

معمارياً، تم تطوير المنصة كواجهة تفاعلية حديثة باستخدام Next.js 16 مع MapLibre GL ونظام المعالجة الخلفية عبر Python FastAPI، وتعمل محلياً على المعالج المركزي (CPU) بزمن تشغيل يبلغ \textbf{1.485 ثانية} وبكلفة سحابية صفرية (\$0)، مما يجعلها منظومة جاهزة لدعم التخطيط الحضري والحوكمة البيئية ومبادرة السعودية الخضراء.

\vspace*{0.8cm}
\noindent
\textbf{الكلمات المفتاحية:} الاستشعار عن بعد، كشف التغير، نماذج الأساس الفضائية، AlphaEarth، Dynamic World، الشبكات العصبية السيامية، Sentinel-2، البيئات شديدة الجفاف، رؤية السعودية 2030.
\newpage
